\documentclass[../numerical_methods.tex]{subfiles}
\begin{document}
\section{Aula 12 - ?? de Maio, 2026}
\subsection{Motivações}
\begin{itemize}
	\item Matrizes Ortonormais;
	\item Método de Cholesky.
\end{itemize}
\subsection{Método de Cholesky}
Apesar de muito útil, o critério para podermos decompor uma matriz em LU ainda é meio rígido; sendo assim, seria interessante termos uma alternativa a ele também! Aqui, entra o método conhecido como \textit{método de Cholesky}, mas antes de vê-lo, temos que definir:
\begin{def*}
	Uma matriz simétrica \(A\in \mathbb{M}_{n}\) é dita \textbf{simétrica positiva definida} (SPD), se, para todo vetor não nulo \(v\in \mathbb{R}^{n}\), temos
	\[
		v \cdot Av > 0.
	\]
\end{def*}
\begin{prop*}
	Cada um dos testes abaixo é uma condição necessária e suficiente para uma matriz simétrica \(A\in \mathbb{M}_{n}\) ser SPD:
	\begin{itemize}
		\item[i)] Para toda menor principal \(A_{k}\), temos \(\det{(A_{k})} > 0,\; k=1, \dotsc, n\);
		\item[ii)] Todos os autovalores de A são positivos.
	\end{itemize}
\end{prop*}
\begin{example}
	A matriz
	\[
		A = \begin{bmatrix}
			3 & 2 & 0 \\
			2 & 5 & 1 \\
			0 & 1 & 3
		\end{bmatrix}
	\]
	é SPD.
\end{example}

Dada uma matriz \(A\in \mathbb{M}_{n}\), queremos \(H\in \mathbb{M}_{n}\) tal que
\[
	A = H \cdot H^{T},
\]
com
\[
	H = \begin{bmatrix}
		h_{11} &        &        &        &        \\
		h_{21} & h_{22} &        &                 \\
		h_{31} & h_{32} & h_{33} &        &        \\
		\vdots & \vdots & \vdots & \ddots &        \\
		h_{n1} & h_{n2} & h_{n3} & \dotsc & h_{nn}
	\end{bmatrix}, \quad h_{ii} > 0,\; i = 1, \dotsc , n.
\]

Para a diagonal, temos
\[
	h_{ii} \Rightarrow a_{ii} = h_{ik}(h_{ki})^{T} = \sum\limits_{k=1}^{i} h_{ik}^{2} = \sum\limits_{k=1}^{i-1}h_{ik}^{2}+h_{ii}^{2},
\]
tal que
\[
	h_{ii} = \biggl(a_{i} - \sum\limits_{k=1}^{i-1}h_{ik}^{2}\biggr)^{\frac{1}{2}}.
\]

Fora da diagonal, com i maior que j,
\[
	a_{ij}=\sum\limits_{k=1}^{j}h_{ik}h_{jk} = \sum\limits_{k=1}^{j-1}h_{ik}h_{jk} + h_{ij}h_{jj},
\]
donde segue que
\[
	h_{ij} = \frac{a_{ij} - \sum\limits_{k=1}^{j-1}h_{ik}h_{jk}}{h_{jj}}.
\]
\begin{crl*}
	Seja \(A\in \mathbb{M}_{n}\) SPD; então, existe uma única matriz triangular inferior \(H\in \mathbb{M}_{n}\) tal que
	\[
		A = H \cdot H^{T}.
	\]
\end{crl*}
Podemos, então, estimar a complexidade para fatorar uma matriz quadrada de tamanho n com os dois métodos que conhecemos até agora:
\begin{itemize}
	\item LU: requer \(\displaystyle \frac{2n^{3}}{3}\) flops;
	\item Cholesky: requer \(\displaystyle \frac{n^{3}}{3}\) flops.
\end{itemize}
Portanto, o método de Cholesky é 50\% mais eficiente que o método LU! Além disso, como ele usa apenas a matriz H (e sua transposta, que é basicamente a mesma coisa), este método tem economia de memória, porque armazena só uma matriz, enquanto o LU armazena a matriz L e a U.

\subsection{Eliminação de Gauss: Termo Geral}
\begin{def*}
	Dois sistemas lineares são \textbf{equivalentes} se tiverem a mesma solução. \(\square\)
\end{def*}

Uma ferramenta extremamente útil que vamos usar é, na verdade, apenas um novo jeito de ver o que já sabemos: as operações elementares. Denotando por \(L_{i}\) a i-ésima linha (equação) de um sistema linear, temos 3 operações elementares:
\begin{itemize}
	\item[1)] Trocar duas linhas no sistema -- \(L_{i} \leftrightarrow L_{j}\);
	\item[2)] Multiplicar uma linha por um escalar não nulo \(\lambda \neq 0\): \(L_{i} \leftrightarrow \lambda L_{i}\);
	\item[3)] Somar a uma linha um múltiplo de uma outra linha: \(L_{i} \leftarrow L_{i} + \lambda L_{j}\).
\end{itemize}

\begin{theorem*}
	Se um sistema linear é obtido a partir de um outro sistema através de uma sequência finita de operações elementares, então eles são equivalentes.
\end{theorem*}

Ao trocar duas linhas no sistema, a ideia do que acontece é \(L_{i}\leftrightarrow L_{j}:\)
\[
	I= \begin{bmatrix}
		1 &        &                                 &        &                                 &        &   \\
		  & \ddots &                                 &        &                                 &        &   \\
		  &        & \underbrace{1}_{\text{linha i}} &        &                                 &        &   \\
		  &        &                                 & \ddots &                                 &        &   \\
		  &        &                                 &        & \underbrace{1}_{\text{linha j}} &        &   \\
		  &        &                                 &        &                                 & \ddots &   \\
		  &        &                                 &        &                                 &        & 1
	\end{bmatrix} \Rightarrow M = \begin{bmatrix}
		1 &        &        &        &                                 &        &   \\
		  & \ddots &        &        &                                 &        &   \\
		  &        & 0      & \dotsc & \underbrace{1}_{\text{linha j}} &        &   \\
		  &        & \vdots & \ddots & \vdots                          &        &   \\
		  &        & 1      & \dotsc & \underbrace{1}_{\text{linha i}} &        &   \\
		  &        &        &        &                                 & \ddots &   \\
		  &        &        &        &                                 &        & 1
	\end{bmatrix}
\]
A multiplicação de uma linha por um escalar \(\lambda \neq 0\) é
\[
	L_{i}\leftarrow \lambda_{ii}L_{i}  \rightsquigarrow M = \begin{bmatrix}
		1 &        &   &              &   &        &   \\
		  & \ddots &   &              &   &        &   \\
		  &        & 1 &              &   &        &   \\
		  &        &   & \lambda_{ii} &   &        &   \\
		  &        &   &              & 1 &        &   \\
		  &        &   &              &   & \ddots &   \\
		  &        &   &              &   &        & 1
	\end{bmatrix}
\]
Somar um múltiplo de uma outra linha a uma linha:
\[
	L_{i}\leftarrow L_{i}+\lambda_{ij}L_{j} \rightsquigarrow M = \begin{bmatrix}
		1 &        &              &        &   &        &   \\
		  & \ddots &              &        &   &        &   \\
		  &        & 1            &        &   &        &   \\
		  &        &              & \ddots &   &        &   \\
		  &        & \lambda_{ij} &        & 1 &        &   \\
		  &        &              &        &   & \ddots &   \\
		  &        &              &        &   &        & 1
	\end{bmatrix} = I + \lambda_{ij}e_{i}e_{j}^{T},
\]
onde I é a matriz identidade.

Algumas propriedades da operação elementar do tipo (3) inclui a preservação do determinante: se
\(A = \begin{bmatrix}
	\dotsc \\
	L_{j}  \\
	\dotsc \\
	L_{i}  \\
	\dotsc
\end{bmatrix}\) e \(\tilde{A} = \begin{bmatrix}
	\dotsc              \\
	L_{j}               \\
	\dotsc              \\
	L_{i}+\lambda L_{j} \\
	\dotsc
\end{bmatrix}\), então
\[
	\det{(\tilde{A})} = \det \begin{bmatrix}
		\dotsc                \\
		L_{j}                 \\
		\dotsc                \\
		L_{i} + \lambda L_{j} \\
		\dotsc
	\end{bmatrix} = \det \begin{bmatrix}
		\dotsc \\
		L_{j}  \\
		\dotsc \\
		L_{i}  \\
		\dotsc
	\end{bmatrix} + \lambda \det \begin{bmatrix}
		\dotsc \\
		L_{j}  \\
		\dotsc \\
		L_{j}  \\
		\dotsc
	\end{bmatrix} = \det{(A)}.
\]

Uma outra é no cálculo da inversa:
\[
	M = \begin{bmatrix}
		1 &        &              &        &   &        &   \\
		  & \ddots &              &        &   &        &   \\
		  &        & 1            &        &   &        &   \\
		  &        &              & \ddots &   &        &   \\
		  &        & \lambda_{ij} &        & 1 &        &   \\
		  &        &              &        &   & \ddots &   \\
		  &        &              &        &   &        & 1
	\end{bmatrix}\Rightarrow M^{-1}=\begin{bmatrix}
		1 &        &               &        &   &        &   \\
		  & \ddots &               &        &   &        &   \\
		  &        & 1             &        &   &        &   \\
		  &        &               & \ddots &   &        &   \\
		  &        & -\lambda_{ij} &        & 1 &        &   \\
		  &        &               &        &   & \ddots &   \\
		  &        &               &        &   &        & 1
	\end{bmatrix}
\]

Dado um sistema linear \(Ax = b\) de ordem n, onde A possui todos os menores principais não singulares; isto é, \(\det{(A_{k})}\neq 0\), onde \(k=1, \dotsc , n\).

Nosso objetivo é obter um sistema linear triangular superior equivalente a \(Ax = b\) usando apenas as operações elementares que preservam determinante, ou seja, operações do tipo 3, conhecidas como \textbf{escalonamento}.
Depois, basta usar o algoritmo de substituições regressivas no sistema escalonado.

Para isso, vamos utilizar a \textbf{matriz aumentada} \([ A | b ]\) e alguns passos:

\textbf{\underline{Passo 1}:} anular todos os elementos abaixo do pivô \(a_{11}^{(1)}\):
\[
	\begin{bmatrix}
		a_{11}^{(1)} & a_{12}^{(1)} & a_{13}^{(1)} & \dotsc & a_{1n}^{(1)} & \bigm| & b_{1}^{(1)} \\
		a_{21}^{(1)} & a_{22}^{(1)} & a_{23}^{(1)} & \dotsc & a_{2n}^{(1)} & \bigm| & b_{2}^{(1)} \\
		a_{31}^{(1)} & a_{32}^{(1)} & a_{33}^{(1)} & \dotsc & a_{3n}^{(1)} & \bigm| & b_{3}^{(1)} \\
		\vdots       & \vdots       & \vdots       & \ddots & \vdots       & \bigm| & \vdots      \\
		a_{n1}^{(1)} & a_{n2}^{(1)} & a_{n3}^{(1)} & \dotsc & a_{nn}^{(1)} & \bigm| & b_{n}^{(1)} \\
	\end{bmatrix}
\]
Para \(i=2, \dotsc , n\), faça:
\[
	\boxed{L_{i}^{(2)}\leftarrow L_{i}^{(1)}+m_{i1}L_{1}^{(1)}, \quad m_{i1} = - \frac{a_{i1}^{(1)}}{a_{11}(1)}}
\]

Ao final do passo 1, apenas a primeira linha estará sem as colunas zeradas, tornando a matriz em:
\[
	\begin{bmatrix}
		a_{11}^{(1)} & a_{12}^{(1)} & a_{13}^{(1)} & \dotsc & a_{1n}^{(1)} & \bigm| & b_{1}^{(1)} \\
		0            & a_{22}^{(2)} & a_{23}^{(2)} & \dotsc & a_{2n}^{(2)} & \bigm| & b_{2}^{(2)} \\
		0            & a_{32}^{(2)} & a_{33}^{(2)} & \dotsc & a_{3n}^{(2)} & \bigm| & b_{3}^{(2)} \\
		\vdots       & \vdots       & \vdots       & \ddots & \vdots       & \bigm| & \vdots      \\
		0            & a_{n2}^{(2)} & a_{n3}^{(2)} & \dotsc & a_{nn}^{(2)} & \bigm| & b_{n}^{(2)} \\
	\end{bmatrix}
\]

\textbf{\underline{Passo 2}:} anular todos os elementos abaixo do pivô \(a_{22}^{(2)}\), fazendo:
\[
	\boxed{L_{i}^{(3)}\leftarrow L_{i}^{(2)} + m_{i2}L_{2}^{(2)}, \quad m_{i2} = -\frac{a_{i2}}{a_{k}}},
\]
resultando, ao final, em
\[
	\begin{bmatrix}
		a_{11}^{(1)} & a_{12}^{(1)} & a_{13}^{(1)} & \dotsc & a_{1n}^{(1)} & \bigm| & b_{1}^{(1)} \\
		0            & a_{22}^{(2)} & a_{23}^{(2)} & \dotsc & a_{2n}^{(2)} & \bigm| & b_{2}^{(2)} \\
		0            & 0            & a_{33}^{(3)} & \dotsc & a_{3n}^{(3)} & \bigm| & b_{3}^{(3)} \\
		\vdots       & \vdots       & \vdots       & \ddots & \vdots       & \bigm| & \vdots      \\
		0            & 0            & a_{n3}^{(3)} & \dotsc & a_{nn}^{(3)} & \bigm| & b_{n}^{(3)} \\
	\end{bmatrix}
\]

\textbf{\underline{Passo n-1}:} anular todos os elementos abaixo do pivô \(a_{n-1, n-1}^{(n-1)}\), utilizando:
\[
	\boxed{L_{i}^{(n)}\leftarrow L_{i}^{(n-1)}+m_{i, n-1}L_{n-1}^{(n-1)},\quad m_{i, n-1} = - \frac{a_{i, n-1}^{(n-1)}}{a_{n-1, n-1}^{(n-1)}}},
\]
obtendo a matriz diagonal superior
\[
	\begin{bmatrix}
		a_{11}^{(1)} & a_{12}^{(1)} & a_{13}^{(1)} & \dotsc & a_{1n}^{(1)} & \bigm| & b_{1}^{(1)} \\
		0            & a_{22}^{(2)} & a_{23}^{(2)} & \dotsc & a_{2n}^{(2)} & \bigm| & b_{2}^{(2)} \\
		0            & 0            & a_{33}^{(3)} & \dotsc & a_{3n}^{(3)} & \bigm| & b_{3}^{(3)} \\
		\vdots       & \vdots       & \vdots       & \ddots & \vdots       & \bigm| & \vdots      \\
		0            & 0            & 0            & \dotsc & a_{nn}^{(n)} & \bigm| & b_{n}^{(n)} \\
	\end{bmatrix}
\]

Portanto, um passo geral \(k=1, \dotsc , n-1\) é realizado com
\[
	\boxed{L_{i}^{(k+1)}\leftarrow L_{i}^{(k)} + m_{ik}L_{k}^{(k)},\quad m_{ik} = - \frac{a_{ik}^{(k)}}{a_{kk}^{(k)}}}
\]
\end{document}
